\section{The extreme decomposition of a Gibbs simplex}

The abstract track of the previous section proves (7.7)(a) at the generality of Remark (7.13).
This section records the theorem for specifications proper, with the parts (b)--(d) and the
decomposition (7.25)--(7.26) with its corollaries.

\begin{theorem}[Georgii, Theorem (7.7)]
    \label{thm:ext-77-concrete}
    \uses{def:spec, def:gibbs-meas, def:ql-tail}
    \lean{MeasureTheory.GibbsMeasure.not_mem_extremePoints_G_of_tail_prob,
      MeasureTheory.GibbsMeasure.tailTrivial_of_mem_extremePoints_G,
      MeasureTheory.GibbsMeasure.mem_extremePoints_G_of_isTailTrivial,
      MeasureTheory.GibbsMeasure.eq_of_forall_measurableSet_tail_eq,
      MeasureTheory.GibbsMeasure.exists_tail_eq_one_eq_zero_of_mem_extremePoints,
      MeasureTheory.GibbsMeasure.mutuallySingular_of_mem_extremePoints}
    \leanok

    Let $\gamma$ be a specification over a countable parameter set.
    \begin{enumerate}
        \item[(a)] $\mu \in \mathcal{G}(\gamma)$ is extreme iff $\mu$ is trivial on $\Tail$.
        \item[(c)] Each $\mu \in \mathcal{G}(\gamma)$ is determined by its restriction to
          $\Tail$: two Gibbs measures agreeing on every tail event are equal.
        \item[(d)] Distinct extreme Gibbs measures $\mu \neq \nu$ are separated by a \emph{tail}
          event $A$ with $\mu(A) = 1$ and $\nu(A) = 0$; in particular they are mutually singular.
    \end{enumerate}
    Part (c) goes through (b): the mixture $\rho = \tfrac12\mu + \tfrac12\nu$ is Gibbs, both
    measures have $\Tail$-measurable densities with respect to it, and agreement on $\Tail$
    forces the densities to agree. No standard Borel hypothesis and no non-emptiness hypothesis
    appears in (c) or (d).
\end{theorem}

\begin{proposition}[The Gibbs kernel, Georgii (7.25)]
    \label{prop:ext-gibbs-kernel}
    \uses{thm:ext-77-concrete}
    \lean{MeasureTheory.GibbsMeasure.gibbsKernel,
      MeasureTheory.GibbsMeasure.gibbsKernel_mem_G,
      MeasureTheory.GibbsMeasure.condExp_ae_eq_gibbsKernel,
      MeasureTheory.GibbsMeasure.isPAKernel_gibbsKernel,
      MeasureTheory.GibbsMeasure.ae_mem_extremePoints_G_gibbsKernel}
    \leanok

    For a standard Borel state space and $\mathcal{G}(\gamma) \neq \emptyset$ there is a
    $(\mathcal{G}(\gamma), \Tail)$-kernel: a tail-measurable kernel $\omega \mapsto
    \pi(\cdot \mid \omega)$ with values in $\mathcal{G}(\gamma)$ which is simultaneously a
    version of $\mu(\cdot \mid \Tail)$ for \emph{every} $\mu \in \mathcal{G}(\gamma)$, and whose
    values are almost surely extreme.
\end{proposition}

\begin{theorem}[The extreme decomposition, Georgii (7.26)]
    \label{thm:ext-726}
    \uses{prop:ext-gibbs-kernel}
    \lean{MeasureTheory.GibbsMeasure.weightOf,
      MeasureTheory.GibbsMeasure.exists_unique_weight_extremePoints,
      MeasureTheory.GibbsMeasure.eq_weightOf_of_join_eq,
      MeasureTheory.GibbsMeasure.weightOf_join,
      MeasureTheory.GibbsMeasure.bijOn_weightOf}
    \leanok

    Every $\mu \in \mathcal{G}(\gamma)$ is the barycentre `Measure.join (weightOf hG μ)` of a
    unique probability weight `weightOf hG μ` concentrated on the extreme points, and
    $\mu \mapsto w_\mu$ is a bijection onto the probability weights on
    $\operatorname{ex}\mathcal{G}(\gamma)$: the Gibbs simplex.
\end{theorem}

\begin{corollary}[Georgii (7.28)--(7.30)]
    \label{cor:ext-728-730}
    \uses{thm:ext-726}
    \lean{MeasureTheory.GibbsMeasure.weightOf_map,
      MeasureTheory.GibbsMeasure.map_eq_self_iff_weightOf_map_eq_self,
      MeasureTheory.GibbsMeasure.measurePreserving_iff_measurePreserving_weightOf,
      MeasureTheory.GibbsMeasure.linearIndependent_of_mem_extremePoints,
      MeasureTheory.GibbsMeasure.setOf_mem_GP_eq_closure_convexCombosLimitGibbs}
    \leanok

    (7.28): the decomposition commutes with symmetries — `μ` is `τ`-invariant iff its weight is.
    (7.29): $\mathcal{G}(\gamma)$ contains $N$ extreme points iff it contains $N$ measures
    linearly independent over $\mathbb{R}_{\geq 0}^\infty$.
    (7.30): for a finite state space and quasilocal $\gamma$, $\mathcal{G}(\gamma)$ is the closed
    convex hull of the limiting Gibbs distributions (Georgii states this for standard Borel $E$;
    the finite-state restriction is recorded in \texttt{formalization.yaml}).
\end{corollary}
